\chapter{Calcul stochastique}\label{ch:iv1:stochastic-calculus}

\begin{ivtrains}
Le mouvement brownien comme outil~: lemme d'Itô, martingales et temps d'arrêt, principe de réflexion et temps d'atteinte, intégrales de la trajectoire, changement de probabilité de Girsanov et Feynman--Kac. La plupart des résultats s'obtiennent de deux façons, et la seconde sert de vérification.
\end{ivtrains}

\begin{question}[$\star$ \iqroles{QR, QT, QD, Risque} \iqfirm{toutes} \iqlevel{junior}]\label{iq:iv1:stochastic-calculus:1}
$W$ est un mouvement brownien standard. Donnez $\operatorname{Cov}(W_s, W_t)$. Calculez la variance de $W_1 + W_2$, la corrélation de $W_1$ et $W_4$, et la probabilité que $W_1$ et $W_2$ soient toutes deux positives.
\end{question}

\begin{question}[$\star$ \iqroles{QR, QT, QD, Risque} \iqfirm{toutes} \iqlevel{junior}]\label{iq:iv1:stochastic-calculus:2}
Lesquelles de ces expressions sont des martingales~: $W_t$, $W_t^2$, $W_t^2 - t$, $tW_t$, $e^{W_t - t/2}$, $W_t^3$, $W_t^3 - 3tW_t$~? Donnez un seul test qui permette de trancher pour toutes.
\end{question}

\begin{question}[$\star$ \iqroles{QR, QT, QD, Risque} \iqfirm{banque, hedge fund} \iqlevel{junior}]\label{iq:iv1:stochastic-calculus:3}
Une action suit $dS = \mu S\,dt + \sigma S\,dW$ avec $\mu = 10~\%$ et $\sigma = 20~\%$ par an. Appliquez le lemme d'Itô à $\ln S$. Partant de 1, quelles sont la moyenne et la médiane après un an~? Quelle est la probabilité de terminer au-dessus de la moyenne~?
\end{question}

\begin{question}[$\star$ \iqroles{QR, QT, QD, Risque} \iqfirm{toutes} \iqlevel{junior}]\label{iq:iv1:stochastic-calculus:4}
Calculez $\mathbb{E}W_t^4$ et $\mathbb{E}e^{aW_t}$ à partir de la loi normale. Que vaut $\mathbb{E}e^{2W_1}$~? Déterminez la corrélation de $W_1^2$ avec $W_4^2$.
\end{question}

\begin{question}[$\star$ \iqroles{QR, QT, QD, Risque} \iqfirm{toutes} \iqlevel{junior}]\label{iq:iv1:stochastic-calculus:5}
Montrez que $cW_{t/c^2}$ est un mouvement brownien pour tout $c \ne 0$. Utilisez ce résultat pour relier $P(W_4 > 2)$ à une probabilité à l'instant 1, et donnez sa valeur. Un rendement a une volatilité quotidienne de 1~\%. Quelle est la volatilité annuelle, avec 252 jours et des rendements indépendants~?
\end{question}

\begin{question}[$\star$ \iqroles{QR, QT, QD, Risque} \iqfirm{toutes} \iqlevel{junior}]\label{iq:iv1:stochastic-calculus:6}
Découpez $[0, 1]$ en 1000 pas égaux. Calculez la moyenne et l'écart type de la somme des accroissements carrés de $W$, et donnez la moyenne de la somme des valeurs absolues des accroissements~? Que devient chacune de ces quantités quand les pas se réduisent~?
\end{question}

\begin{question}[$\star\star$ \iqroles{QR, QT, QD} \iqfirm{banque, hedge fund} \iqlevel{junior, confirmé}]\label{iq:iv1:stochastic-calculus:7}
Calculez $\int_0^TW_t\,dW_t$ à l'aide du lemme d'Itô. Donnez sa moyenne et sa variance. On l'approche par une somme sur de petits pas, en évaluant $W$ à l'extrémité gauche de chaque pas, puis à l'extrémité droite, puis au milieu. Vers quoi converge chacune de ces sommes~?
\end{question}

\begin{question}[$\star\star$ \iqroles{QR, QT, QD} \iqfirm{banque, hedge fund} \iqlevel{junior, confirmé}]\label{iq:iv1:stochastic-calculus:8}
Soit $I_T = \int_0^TW_t\,dt$. Déterminez sa loi, par une intégration par parties, et vérifiez la variance par une intégrale double. Quelle est sa corrélation avec $W_T$~? Évaluez la variance pour $T = 3$.
\end{question}

\begin{question}[$\star\star$ \iqroles{QR, QT, QD, Risque} \iqfirm{banque, hedge fund} \iqlevel{junior, confirmé}]\label{iq:iv1:stochastic-calculus:9}
Énoncez le principe de réflexion, et déduisez la loi de $M_T = \max_{s \le T}W_s$. Quelles sont $P(M_1 \ge 1)$ et $\mathbb{E}M_1$~? Comment votre réponse change-t-elle si la barrière n'est contrôlée qu'aux clôtures quotidiennes, avec 252 clôtures dans l'année~?
\end{question}

\begin{question}[$\star\star$ \iqroles{QR, QT, QD, Risque} \iqfirm{toutes} \iqlevel{junior, confirmé}]\label{iq:iv1:stochastic-calculus:10}
$T_a$ est le premier instant où $W$ atteint $a > 0$. Montrez que $T_a$ est fini presque sûrement, et calculez sa médiane pour $a = 1$. Quelle est $P(T_1 > 100)$~? Pourquoi son espérance est-elle infinie~?
\end{question}

\begin{question}[$\star\star$ \iqroles{QR, QT, QD, Risque} \iqfirm{toutes} \iqlevel{junior}]\label{iq:iv1:stochastic-calculus:11}
$W$ part de 0 et s'arrête dès qu'elle atteint $-1$ ou $2$. À l'aide de deux martingales, calculez la probabilité qu'elle finisse en 2 et l'espérance du temps d'arrêt. Quelle est la formule générale pour $-a$ et $b$~? Que se passe-t-il quand $a \to \infty$~?
\end{question}

\begin{question}[$\star\star$ \iqroles{QR, QT, QD, Risque} \iqfirm{hedge fund, banque} \iqlevel{junior, confirmé}]\label{iq:iv1:stochastic-calculus:12}
$X_t = \mu t + \sigma W_t$ avec $\mu > 0$ et $X_0 = 0$. Déterminez une martingale exponentielle, et utilisez-la pour établir la probabilité que $X$ tombe un jour à $-a$. Une stratégie a une moyenne annuelle de 10~\% et une volatilité de 20~\%, sans capitalisation. Quelle est la probabilité qu'elle perde un jour 20~\% depuis le départ~?
\end{question}

\begin{question}[$\star\star$ \iqroles{QR, QT, QD, Risque} \iqfirm{banque, hedge fund} \iqlevel{junior, confirmé}]\label{iq:iv1:stochastic-calculus:13}
Pour $dS = \mu S\,dt + \sigma S\,dW$, appliquez le lemme d'Itô à $S^p$. Calculez $\mathbb{E}S_T^2$, et trouvez le $p$ qui fait de $S^p$ une martingale, avec $\mu = 10~\%$ et $\sigma = 20~\%$. Partant de 1, quelle est la probabilité d'atteindre 2 avant $\frac12$~?
\end{question}

\begin{question}[$\star\star$ \iqroles{QR, QT, QD, Risque} \iqfirm{banque, hedge fund} \iqlevel{junior, confirmé}]\label{iq:iv1:stochastic-calculus:14}
Sachant $W_1 = x$, quelles sont l'espérance et la variance conditionnelles de $W_s$, $0 < s < 1$~? Sachant $W_1 = 0$, quelle est la probabilité que $W$ atteigne le niveau 1 pendant l'intervalle~? Établissez la formule par réflexion, pour tout $x$.
\end{question}

\begin{question}[$\star\star$ \iqroles{QR, QT, QD, Risque} \iqfirm{toutes} \iqlevel{junior, confirmé}]\label{iq:iv1:stochastic-calculus:15}
Résolvez l'équation d'Ornstein--Uhlenbeck $dX = -\kappa X\,dt + \sigma\,dW$. Donnez l'espérance et la variance à $t$, la variance stationnaire et la demi-vie d'un choc. Évaluez-les pour $\kappa = 2$, $\sigma = 1$, $t = 1$, $X_0 = 3$.
\end{question}

\begin{question}[$\star\star$ \iqroles{QR, QT, QD, Risque} \iqfirm{banque, hedge fund} \iqlevel{junior}]\label{iq:iv1:stochastic-calculus:16}
$W^1$ et $W^2$ ont une corrélation $\rho$. Exprimez l'un en fonction de l'autre et d'un mouvement indépendant. Utilisez le lemme d'Itô pour calculer $\mathbb{E}W^1_tW^2_t$. Pour $\rho = 0.5$ et $t = 1$, calculez $\operatorname{Var}(W^1_t + W^2_t)$ et $\mathbb{E}[(W^1_t)^2(W^2_t)^2]$.
\end{question}

\begin{question}[$\star\star$ \iqroles{QR, QT, QD} \iqfirm{banque, hedge fund} \iqlevel{junior, confirmé}]\label{iq:iv1:stochastic-calculus:17}
Montrez directement, en complétant le carré, que $\mathbb{E}[f(W_T + \mu T)] = \mathbb{E}[f(W_T)e^{\mu W_T - \mu^2T/2}]$. Utilisez cette identité pour estimer $P(W_1 > 4)$ par simulation, avec une dérive déformée. De combien cela réduit-il l'erreur type, pour un million de trajectoires~?
\end{question}

\begin{question}[$\star\star$ \iqroles{QR, QT, QD, Risque} \iqfirm{banque, hedge fund} \iqlevel{junior}]\label{iq:iv1:stochastic-calculus:18}
Une action a $\mu = 10~\%$, $\sigma = 20~\%$, $S_0 = 100$, et le taux est nul. Un digital à un an paie 1 si $S_1 > 100$. Quelles sont la probabilité historique, la probabilité risque-neutre et le prix~? Quel est le prix de marché du risque, et comment relie-t-il les deux mesures~? Le prix dépend-il de $\mu$~?
\end{question}

\begin{question}[$\star\star$ \iqroles{QR, QT, QD} \iqfirm{banque, hedge fund} \iqlevel{junior, confirmé}]\label{iq:iv1:stochastic-calculus:19}
Énoncez la formule de Feynman--Kac pour $u_t + \mu u_x + \frac12\sigma^2u_{xx} - ru = 0$, avec la condition terminale $u(T, x) = g(x)$. Résolvez l'équation avec $\mu = r = 0$, $\sigma = 1$, d'abord pour $g = x^2$, puis pour $g = e^x$. Vérifiez chaque résultat en le substituant dans l'équation. Comment un taux $r$ les modifie-t-il~?
\end{question}

\begin{question}[$\star\star$ \iqroles{QR, QT, QD, Risque} \iqfirm{hedge fund} \iqlevel{junior, confirmé}]\label{iq:iv1:stochastic-calculus:20}
Soit $X_t = e^{W_t - t/2}$. Quelle équation résout-il~? Quelles sont $\mathbb{E}X_T$, la médiane et $P(X_T < 1)$ pour $T = 1$ et $T = 100$~? Quel est le comportement à long terme d'une trajectoire unique~? Comment concilier ce résultat avec la moyenne~?
\end{question}

\begin{question}[$\star\star\star$ \iqroles{QR, QT, QD} \iqfirm{banque, hedge fund} \iqlevel{confirmé}]\label{iq:iv1:stochastic-calculus:21}
Quelles sont l'espérance et la variance conditionnelles de $\int_0^TW_t\,dt$ sachant $W_T$~? Évaluez-les pour $T = 1$. Quelle part de la variance de l'intégrale est expliquée par la valeur finale~?
\end{question}

\begin{question}[$\star\star\star$ \iqroles{QR, QT} \iqfirm{hedge fund} \iqlevel{junior, confirmé}]\label{iq:iv1:stochastic-calculus:22}
Soit $A$ la fraction de $[0, 1]$ pendant laquelle $W$ reste au-dessus de 0. Énoncez la loi de $A$. Donnez la moyenne, la variance et $P(A \le 0.1)$. Le profit cumulé d'un trader se comporte comme $W$. Qu'est-ce que cela dit de la fréquence à laquelle un trader sans talent est en avance pendant au moins 90~\% de l'année~?
\end{question}

\begin{question}[$\star\star\star$ \iqroles{QR, QT} \iqfirm{hedge fund} \iqlevel{junior, confirmé}]\label{iq:iv1:stochastic-calculus:23}
Soit $M_T$ le maximum courant de $W$. À partir du principe de réflexion, calculez $P(M_T \ge a, W_T \le x)$ pour $x \le a$. Montrez que $M_T - W_T$ a la même loi que $|W_T|$. Une stratégie sans dérive a une volatilité de 20~\% sans capitalisation. Quelle est la distance moyenne sous son plus haut après un an, et quelle est la probabilité que cet écart dépasse 20~\%~?
\end{question}

\begin{question}[$\star\star\star$ \iqroles{QR, QT, QD, Risque} \iqfirm{hedge fund, banque} \iqlevel{confirmé}]\label{iq:iv1:stochastic-calculus:24}
$X_t = \mu t + \sigma W_t$ avec $\mu > 0$, et $\tau$ est le premier instant où $X$ atteint $a > 0$. Utilisez les martingales pour établir $\mathbb{E}\tau$ et $\operatorname{Var}\tau$. Évaluez-les pour $a = 1$, $\mu = 0.1$ et $\sigma = 0.2$. Quelle est $P(\tau > 20)$, et de quelle loi s'agit-il~?
\end{question}

\begin{question}[$\star\star\star$ \iqroles{QR, QT} \iqfirm{hedge fund} \iqlevel{confirmé}]\label{iq:iv1:stochastic-calculus:25}
Utilisez une martingale exponentielle et le théorème d'arrêt pour montrer que $\mathbb{E}e^{-\lambda T_a} = e^{-a\sqrt{2\lambda}}$, où $T_a$ est le premier instant où $W$ atteint $a > 0$. Pourquoi faut-il choisir la racine positive~? Que découvre-t-on quand $\lambda \to 0$, et que dit la dérivée en ce point~? Évaluez pour $a = 1$ et $\lambda = \frac12$.
\end{question}

\begin{question}[$\star\star\star$ \iqroles{QR, QT, QD, Risque} \iqfirm{hedge fund, banque} \iqlevel{confirmé}]\label{iq:iv1:stochastic-calculus:26}
Un actif suit un mouvement brownien géométrique de volatilité 20~\%. Partant de 1, calculez le temps moyen pour atteindre 2 si $\mu = 10~\%$. Si $\mu = 1~\%$, quelle est la probabilité qu'il atteigne un jour 2~? Que se passe-t-il pour $\mu = 2~\%$~? Vérifiez le second résultat à l'aide d'une martingale de la forme $S^p$.
\end{question}

\begin{question}[$\star\star\star$ \iqroles{QR, QT, QD} \iqfirm{hedge fund, banque} \iqlevel{confirmé}]\label{iq:iv1:stochastic-calculus:27}
Le lemme d'Itô ne s'applique pas à $|x|$, qui présente un point anguleux en 0. En utilisant la formule de Tanaka, écrivez $|W_t|$ comme une intégrale stochastique plus un temps local $L_t$. Calculez $\mathbb{E}L_1$. Que mesure $L$, et comment pourrait-on l'estimer à partir d'une trajectoire simulée~? Quel autre processus a la même loi que $L_t$~?
\end{question}

\begin{question}[$\star\star\star$ \iqroles{QR, QT} \iqfirm{hedge fund} \iqlevel{confirmé}]\label{iq:iv1:stochastic-calculus:28}
Calculez la moyenne et la variance de $Q = \int_0^1W_t^2\,dt$ par des intégrales doubles. Vérifiez ensuite les deux résultats à l'aide du développement $W_t = \sum_k\sqrt{2}\,Z_k\sin((k - \frac12)\pi t)/((k - \frac12)\pi)$, avec des variables normales standard indépendantes $Z_k$, en utilisant $\sum_k((k - \frac12)\pi)^{-2} = \frac12$. Quelle somme faut-il utiliser pour la variance~?
\end{question}

\begin{question}[$\star\star\star\star$ \iqroles{QR, QT} \iqfirm{hedge fund} \iqlevel{confirmé}]\label{iq:iv1:stochastic-calculus:29}
Soit $Q = \int_0^1W_t^2\,dt$. Établissez $\mathbb{E}e^{-\theta^2Q/2} = 1/\sqrt{\cosh\theta}$ à l'aide de la formule de Feynman--Kac, avec l'hypothèse $u(t, x) = \exp(-\frac12A(t)x^2 + B(t))$. Résolvez pour $A$ et $B$. Vérifiez la limite pour $\theta$ petit par rapport à la moyenne de $Q$, qui vaut $\frac12$. Évaluez-la pour $\theta = 1$.
\end{question}

\begin{question}[$\star\star\star\star$ \iqroles{QR, QT} \iqfirm{hedge fund} \iqlevel{confirmé}]\label{iq:iv1:stochastic-calculus:30}
$\tau$ est le premier instant où $W$ quitte $(-a, a)$. En utilisant $\cosh(\theta W_t)e^{-\theta^2t/2}$, calculez $\mathbb{E}e^{-\lambda\tau}$. Développez-la pour obtenir $\mathbb{E}\tau$ et $\operatorname{Var}\tau$. Vérifiez ces résultats à l'aide de la martingale polynomiale $W_t^4 - 6tW_t^2 + 3t^2$. Utilisez ensuite les pôles de la transformée pour donner la vitesse de décroissance de $P(\tau > t)$.
\end{question}
